\subsection{Exakte gemeinsame
Reduktion}\label{exakte-gemeinsame-reduktion}

Der ursprüngliche kausale Paarversuch bleibt vollständig im Raum

\[\adjustbox{max width=\linewidth}{$\displaystyle 
\mathcal K_3=\operatorname{span}\{|p_0\rangle,|R_7\rangle,|b_0\rangle\},
\quad |p_0\rangle=|4,57\rangle,\quad
|R_7\rangle=\frac1{\sqrt7}\sum_{q=1}^7s_q|p_q\rangle.
$}\]

Die acht Paare und ihre Vorzeichen stammen unverändert aus Zeile 0 des
eingefrorenen W; sie stehen im zugehörigen \codeword{REPORT.md}. Die
Isometrie J hat diese drei Vektoren als Spalten. Direkt aus W folgt

\[\adjustbox{max width=\linewidth}{$\displaystyle 
HJ=Jh_3,\qquad Z_4J=Jz_3,\qquad N_bJ=Jb_3,
$}\] \[\adjustbox{max width=\linewidth}{$\displaystyle 
h_3=\begin{pmatrix}0&0&-g\\0&0&\sqrt7g\\-g&\sqrt7g&\Delta\end{pmatrix},
\quad z_3=\operatorname{diag}(-1,1,1),\quad
b_3=\operatorname{diag}(0,0,1).
$}\]

Alle drei Intertwining-Identitäten sind exakt und haben keinen Restterm.
Somit haben beliebige Folgen aus freier H-Entwicklung und Z4 denselben
nativen dreidimensionalen Träger. Hier gilt immer N=2.

Die Empfängerbesetzung besitzt den exakten komprimierten \textbf{Effekt}

\[\adjustbox{max width=\linewidth}{$\displaystyle 
e_3=J^\dagger n_5J=\operatorname{diag}(0,1/7,0).
$}\]

Damit reproduziert der Dreizustandsraum ohne Näherung beide kausalen
Wahrscheinlichkeiten aus \codeword{REPORT.md}: x(4−x)/64 ohne mittleren
Impuls und 9x²/1024 mit Impuls. In beiden Armen sind am dort gewählten
Endzeitpunkt Nf=2 und Nb=0.

\subsection{Wesentliche Grenze: ein terminaler Effekt ist kein
geschlossenes
Instrument}\label{wesentliche-grenze-ein-terminaler-effekt-ist-kein-geschlossenes-instrument}

n5 selbst erhält diesen Dreizustandsraum nicht. Es gilt exakt

\[\adjustbox{max width=\linewidth}{$\displaystyle 
(n_5J-Je_3)^\dagger(n_5J-Je_3)=\operatorname{diag}(0,6/49,0).
$}\]

Eine tatsächliche nichtselektive projektive n5-Messung lässt vom Zustand
R7 das Gewicht \textbf{12/49} außerhalb von K3 zurück. Deshalb gilt die
gemeinsame Reduktion für Entwicklungen und Impulse \textbf{bis zur
abschließenden Messung}. Eine weitere Verwendung desselben Exemplars
nach diesem Instrument wird nicht behauptet. Die Tomografie verwendet
jeweils frische Exemplare desselben präparierten Zustands und genau eine
abschließende Besetzungsmessung.

Ein tatsächlich grobes Lüders-Instrument für den Projektor auf die
gesamten sieben anderen Paarblätter würde K3 erhalten; sein Effekt ist
in K3 diag(0,1,0). Das ist eine zusätzliche mögliche
Instrumentenanforderung. Ein feines Auslesen der sieben Blattlabels und
anschließendes Wegwerfen der Labels ist ein anderes Instrument: Es lässt
aus R7 das Gewicht 6/7 aus K3 heraustreten. Keine dieser
Implementierungen wird als nativ verfügbar vorausgesetzt, sofern sie
nicht ausdrücklich gewährt wird.

\subsection{Vollständige terminale Tomografie auf demselben
Träger}\label{vollstuxe4ndige-terminale-tomografie-auf-demselben-truxe4ger}

Definiere für hermitesche Operatoren

\[\adjustbox{max width=\linewidth}{$\displaystyle 
\mathcal L(O)=i[h_3,O],\qquad\mathcal C(O)=z_3Oz_3.
$}\]

Die folgende Liste besitzt neun linear unabhängige hermitesche
Operatoren:

\[\adjustbox{max width=\linewidth}{$\displaystyle 
\boxed{e_3,b_3,\mathcal Le_3,\mathcal Lb_3,
\mathcal L^2e_3,\mathcal L^2b_3,\mathcal L^3e_3,
\mathcal C\mathcal L^2e_3,\mathcal C\mathcal L^2b_3.}
$}\]

Schreibt man O als reellen Koordinatenvektor

\[\adjustbox{max width=\linewidth}{$\displaystyle 
(O_{00},O_{11},O_{22},\Re O_{01},\Re O_{02},\Re O_{12},
\Im O_{01},\Im O_{02},\Im O_{12}),
$}\]

hat die Matrix dieser neun Spalten die Determinante

\[\adjustbox{max width=\linewidth}{$\displaystyle 
\boxed{\det M=\frac{8\Delta^3g^{10}}{343}\ne0}
\qquad(\Delta g\ne0).
$}\]

Damit spannt der Orbit aus zeitentwickelten terminalen Effekten und
lokaler Impulskonjugation \textbf{Herm(3)} vollständig auf. Der bereits
bewiesene kausale Zeuge und diese Schattenrekonstruktion benutzen
denselben ursprünglichen Tensor, Hamiltonoperator, Träger, Phasenimpuls
und dieselben Besetzungsobservablen. Es handelt sich um
Zustandstomografie bei bekanntem H und bekannten Instrumenten, nicht um
eine gleichzeitige unbekannte Hamilton- und Instrumentenkalibrierung.

Die Operatorwörter sind durch Ableitungen von wirklichen
Endwahrscheinlichkeiten zugänglich:
\(\partial_t^k p_e(0)=\operatorname{tr}(\rho\mathcal L^ke_3)\). Für die
beiden letzten Wörter wird Z4 \textbf{vor} der freien Entwicklung
angewandt;
\(p_e^Z(t)=\operatorname{tr}(\rho z_3e^{ith_3}e_3e^{-ith_3}z_3)\). Die
zugehörige zweite Ableitung liefert das verlangte Wort. Es genügt also
der Zugriff auf die Kurven von n5 und Nb mit und ohne anfänglichen
Impuls, mit Ableitungen bis zur Ordnung drei. Die vollständige
Bestimmung dieser Ableitungen wird als Messdatenanforderung benannt; ein
optimiertes endliches Abtast- und Schätzverfahren wird hier nicht
behauptet.

\subsubsection{Der Impuls ist für normierte Zustandstomografie nicht
zwingend}\label{der-impuls-ist-fuxfcr-normierte-zustandstomografie-nicht-zwingend}

Die rein autonome Zeitfamilie aus e3 und b3 besitzt Dimension acht. Ihre
einzige fehlende hermitesche Richtung kann durch

\[\adjustbox{max width=\linewidth}{$\displaystyle 
Q=\begin{pmatrix}6\Delta&\sqrt7\Delta&-g\\
\sqrt7\Delta&0&\sqrt7g\\-g&\sqrt7g&0\end{pmatrix}
$}\]

dargestellt werden: \([h_3,Q]=0\), \(\operatorname{tr}(e_3Q)=
\operatorname{tr}(b_3Q)=0\), aber \(\operatorname{tr}Q=6\Delta\ne0\).
Für die Differenz zweier normierter Zustände ist diese Richtung daher
ausgeschlossen. Tatsächlich bilden

\[\adjustbox{max width=\linewidth}{$\displaystyle 
I,e_3,b_3,\mathcal Le_3,\mathcal Lb_3,\mathcal L^2e_3,
\mathcal L^2b_3,\mathcal L^3e_3,\mathcal L^4e_3
$}\]

ebenfalls eine Basis; ihre Determinante ist 6Δ³g¹⁰/343. Der bekannte
Wert Trρ=1 ersetzt somit eine weitere Messrichtung. Der Impuls erlaubt
hier eine konkrete Basis mit Ableitungen nur bis Ordnung drei; aus Rang
acht allein folgt kein Tomografiehindernis für normierte Zustände.

\subsection{Bedingte Fehlergrenze}\label{bedingte-fehlergrenze}

Verwende dimensionslose Zeit u=Δt und den Prüfpunkt g/Δ=1/20. Sei y der
Vektor der neun oben angegebenen Operatorerwartungswerte und δy sein
Fehler. Der exakte lineare Inversenrechner liefert mit der
Frobenius-Normschranke

\[\adjustbox{max width=\linewidth}{$\displaystyle 
\|\widehat\rho-\rho\|_{\mathrm{HS}}
\le\frac{\sqrt{9818033}}2\|\delta y\|_2,
$}\] \[\adjustbox{max width=\linewidth}{$\displaystyle 
\frac12\|\widehat\rho-\rho\|_1
\le\frac{\sqrt{29454099}}4\|\delta y\|_2.
$}\]

Sind alle neun Dateneinträge mit Fehler höchstens ε bekannt, gilt für
die zweite Schranke 3√29454099·ε/4. Diese konservative Schranke belegt
Stabilität für die ausdrücklich gegebenen Daten, aber auch eine mögliche
starke Fehlerverstärkung bei schwacher Kopplung. Sie enthält
\textbf{keine} kostenlose Gewinnung genauer Zeitableitungen aus endlich
vielen verrauschten Messungen. Messdauer, Anzahl frischer Exemplare,
Ableitungsschätzung und mögliche Verbesserungen der Messauswahl bleiben
eigene Ressourcenfragen.

\subsection{Beweisumfang}\label{beweisumfang}

\codeword{verify\_common3.py} prüft den eingefrorenen W, die Isometrie,
alle Intertwining-Identitäten, den ausdrücklichen n5-Leckterm, die zwei
kausalen Endwahrscheinlichkeiten, beide symbolischen Determinanten und
die rationale Inversen-Normschranke. \codeword{common3.json} und
\codeword{common3\_OO.json} sind die normalen und optimierten Replays.
Der frühere Bericht und Prüfer wurden nicht geändert.

Geschlossen ist eine \textbf{gemeinsame bedingte Ausführung}:
Zustandsrekonstruktion und kausaler Eingriff im gleichen nativen
N=2-Träger. Offen bleiben die native Auswahl und Präparation der
Zustände, ausführbare Modenkontrolle, kalibrierte terminale Instrumente
und die Ressourcen ihrer Wiederholung. Es folgt keine vollständige
Instrumentenalgebra nach Messung, keine
N=64-Grundzustands-/Polidentifikation und keine räumliche
Interpretation.

\begin{center}\rule{0.5\linewidth}{0.5pt}\end{center}

\section{Nachtragsprüfung A: neue fünfte Norm und richtige
Antwortzustände}\label{nachtragspruxfcfung-a-neue-fuxfcnfte-norm-und-richtige-antwortzustuxe4nde}

\section{Nachtrag A: Was am neuen v1.6.3-Erklärtext wirklich neu
ist}\label{nachtrag-a-was-am-neuen-v1.6.3-erkluxe4rtext-wirklich-neu-ist}

\subsection{Urteil}\label{urteil}

\textbf{Ja, ein belastbarer neuer Rechenschritt ist enthalten:} Die
zuvor fehlende Norm ν5 liegt jetzt mit passendem Quellstand vor. Wir
haben ν1 bis ν5 frisch durch die ursprüngliche exakte
Spurnetzwerkrechnung reproduziert und daraus unabhängig strengere
rationale Ritzgrenzen gewonnen. Die übrigen großen Aussagen des Textes
sind teils bereits bekannt, teils durch v1.6.7 korrigiert. Besonders
problematisch ist die Vermischung verschiedener Ritz-Zustände zu einem
vermeintlich vollständig bestimmten Grundzustand.

Die gesamte neue Anlage wurde gelesen. Sie stimmt mit der derzeitigen
\codeword{EINFACH\_ERKLAERT.md} im Ordner
\codeword{/Users/stefanhamann/Projekte/tfpt-theoryv4/experiments/theory-contracts/universalraum-native-operations-ground-response-20260915/}
überein. Alle für den neuen Nachweis verwendeten Programme und
Ergebnisse sind separat unter \codeword{late\_a/inputs/} eingefroren.
Die vorhandenen v1.6.7- und Mechanismus-Snapshots wurden nicht
verändert.

\subsection{1. ν5 ist jetzt ein Ergebnis, nicht mehr nur eine
Programmeinstellung}\label{ux3bd5-ist-jetzt-ein-ergebnis-nicht-mehr-nur-eine-programmeinstellung}

Der v1.6.7-Snapshot hatte zu Recht festgehalten, dass damals
\codeword{norms\_by\_traces.json} und ein abgeschlossener zugehöriger
Replay fehlten. Dieser konkrete Quellenstand hat sich geändert. Der neue
Checkerhash
\codeword{be72c728e128ae0ef605980a5bae88ea31049e29e2fd17c43167fbd6ad1d0c89}
passt jetzt zum Ergebnis; auch der Common-Hash stimmt. Die vorhandenen
normalen und optimierten Normausgaben sind byteidentisch. Der
gespeicherte JSON-Inhalt ist identisch; nur ein abschließender
Zeilenumbruch unterscheidet die Druckausgabe von der geschriebenen
Datei.

Zusätzlich zu dieser Quellenprüfung wurde die wirkliche Rechnung in
unserem isolierten Paket normal und mit \codeword{-OO} frisch
ausgeführt. Geändert wurde nur der Tensorpfad auf das eingefrorene
Originalarchiv; der eigentliche Spurnetzwerkcode blieb unverändert.
Ergebnis:

\[\adjustbox{max width=\linewidth}{$\displaystyle 
(\nu_0,\nu_1,\nu_2,\nu_3,\nu_4,\nu_5)=
(1,480,439680,575078400,952296652800,1866738327552000).
$}\]

Die fünfte Norm benötigt 840 Wick-Netzwerke, die auf 93 Klassen unter
vorzeichenfreien Graphsymmetrien reduziert werden. Jede tatsächlich
ausgewertete Kontraktion benutzt dünnbesetzte Dictionaries und beliebig
große Python-Ganzzahlen. Hier wurde keine Float-Kontraktion oder Rundung
als exakte fünfte Norm ausgegeben.

Die Methode hat einen nachvollziehbaren Ursprung: Für die
antisymmetrische Paarmatrix M(z) liefert die fermionische
Exponentialnorm \(\det(I+M(z)^\dagger M(z))^{1/2}\). Nach
Logarithmusentwicklung und bosonischer Gaußintegration entstehen die im
Quellcode benutzten Spurzyklen, rationalen Faktoren und
Wick-Permutationen. Die signierte Φ-Kontraktion ist deshalb eine
Normberechnung desselben Paartensors.

\textbf{Reproduktionsgrenze:} Das ist ein frischer Replay der geprüften
ursprünglichen Methode, keine zweite unabhängig neu erfundene Methode
für ν5. Die früheren ν1-ν4-Werte liefern unabhängige bekannte
Vergleichspunkte. Die Quelle vergleicht außerdem 39 kanonische Netzwerke
nur bis n=4 mit ihrem anderen dichten Pfad; daraus folgt kein zweiter
n=5-Beweis. Der Hilfsname \codeword{brute\_force\_nu2} ist irreführend:
Sein Aufbau lässt explizite CAR-Umordnungszeichen aus. Dieser Helfer
wurde nicht als tragender unabhängiger Normnachweis benutzt. Die
signierten Spurnetzwerke bleiben davon getrennt.

Das Gesamt-Replaymanifest der fremden Suite stand beim Einfrieren
weiterhin auf \textbf{RUNNING} mit leerer Laufübersicht. Wir erklären
damit nicht die ganze fremde Suite für abgeschlossen. Unser isolierter
Norm-/Ritznachweis ist hingegen abgeschlossen und reproduziert.

\subsection{2. −1,138 ist eine echte bessere
Variationsgrenze}\label{ist-eine-echte-bessere-variationsgrenze}

Die Unterräume aus den normierten bosonweise geordneten Vektoren
\(v_k=T_+^kF\) sind nicht invariant, aber sie sind gültige
Variationsräume. Da unterschiedliche k orthogonal sind und H die
Bosonzahl höchstens um eins ändert, besitzt die Kompression auf
v0,\ldots,v5 exakt

\[\adjustbox{max width=\linewidth}{$\displaystyle 
(H_6)_{kk}=k\Delta,\qquad
(H_6)_{k,k+1}=g\sqrt{\nu_{k+1}/\nu_k}.
$}\]

Aus der jetzt frisch geprüften ν5 wurde die charakteristische Gleichung
dieser \textbf{sechsdimensionalen reinen Leiterkompression} neu mit
rationaler Arithmetik aufgebaut. Eine unabhängige Sturm-Zählung beweist
bei g/Δ=1/20:

\[\adjustbox{max width=\linewidth}{$\displaystyle 
\boxed{-1.13847420<E_{6}/\Delta<-1.13847419.}
$}\]

Die optional zusätzlich enthaltene Richtung w2 ergibt einen
siebendimensionalen Variationsraum. Ihre Norm ist weiterhin
\(5001523200/229\). Ihre einzige Kopplung innerhalb dieses Raums führt
zu v3:

\[\adjustbox{max width=\linewidth}{$\displaystyle 
\langle\widehat w_2,H\widehat v_3\rangle
=g\sqrt{\|w_2\|^2/\nu_3}.
$}\]

Die v1-Kopplung verschwindet durch w2⊥v2; alle anderen verschwinden
schon wegen der Bosonzahl. Der diagonale Wert ist 2Δ. Damit ist auch
diese Kompression vollständig bestimmt, ohne eine Schließung zu
behaupten. Ihre unabhängige rationale Wurzeleinschließung lautet

\[\adjustbox{max width=\linewidth}{$\displaystyle 
\boxed{-1.13847610<E_{7}/\Delta<-1.13847609.}
$}\]

Nach Rayleigh-Ritz folgt die neue strenge Obergrenze

\[\adjustbox{max width=\linewidth}{$\displaystyle 
\boxed{E_0<-1.13847609\Delta.}
$}\]

Die frühere Grundzustandsuntergrenze −1,158089Δ bleibt ein ausdrücklich
übernommener älterer exakter Eingang. Ebenso wird der ältere N=63-Floor
\(E_h>-1.121899\Delta\) übernommen. Zusammen ergibt sich neu

\[\adjustbox{max width=\linewidth}{$\displaystyle 
\boxed{\epsilon=E_h-E_0>0.01657709\Delta.}
$}\]

Schon die reine Sechserkompression liefert ε\textgreater0,01657519Δ. Der
Siebenerwert ist etwas stärker; beide verbessern den v1.6.7-Floor
0,00773911Δ erheblich. Der ältere N=63-Beweis und das große
Grundzustandsarchiv wurden dabei nicht erneut ausgeführt. Die neue
Folgerung hängt ausdrücklich von diesem beibehaltenen exakten Floor ab.

\textbf{Entscheidende Unterscheidung:} H6/H7 sind keine sechste oder
siebte Fortsetzung der in v1.6.7 berechneten vollständigen
H-Lanczos-Kette. Schon deren fünfter Vektor enthält v4+w2. Die
Leiterkompression ist dennoch eine gültige Ritzrechnung und darf deshalb
ihre Energieobergrenze verbessern. Ein kleines lokales
Nichtschlussverhältnis oder eine flacher werdende Ritzfolge beweist
keine Konvergenz zum wahren Grundzustand. Die Behauptung „nur sehr nahe
an mit fünf Zahlen lösbar`` bleibt unbegründet.

